\documentclass[11pt]{article}
\usepackage[margin=1in]{geometry}
\usepackage{amsmath,amssymb}
\newcommand{\finalanswer}[1]{\par\medskip\noindent\fbox{\parbox{0.94\linewidth}{\textbf{Answer:} #1}}}
\title{STAT 412 Homework 5 --- Question 18}
\author{}\date{}
\begin{document}\maketitle
\section*{Problem}
The government limit is $7260$ ppm. A process is sampled $25$ times; historical
standard deviation is $\sigma=140$ ppm. Part (a) asks for the probability that
the sample average would exceed the limit if the population mean equals the
limit.
Part (b) asks for $P(\bar X\ge7344\mid \mu=7260)$ and the conclusion.

\section*{Work}
If $\mu=7260$, then
\[
\mu_{\bar X}=7260,\qquad
\sigma_{\bar X}=\frac{140}{\sqrt{25}}=28.
\]
The threshold is exactly the mean:
\[
z=\frac{7260-7260}{28}=0.
\]
Therefore
\[
P(\bar X>7260)=P(Z>0)=0.5000.
\]
For the observed sample mean,
\[
z=\frac{7344-7260}{140/\sqrt{25}}=\frac{84}{28}=3.00.
\]
Therefore
\[
P(\bar X\ge7344\mid \mu=7260)=P(Z\ge3.00)=1-0.9987=0.0013.
\]
This probability is negligible, so the observed $\bar x=7344$ is evidence that
the population mean exceeds the government limit.

\finalanswer{(a) $0.5000$.
(b) $0.0013$; select \textbf{is} negligible and $\bar x$ \textbf{is} evidence.}
\end{document}
